% Ampliación para insertar después de hibridacion.tex y del módulo térmico.
% La ecuación electrónica completa y las torres permanecen en sus secciones.
\section{Historique électronique et conservation de la masse sous transport}
\label{md:r27:electron}

L'équation \eqref{md:compos:electron-completo} évalue une section centrale persistante et conserve son retour d'action. La trajectoire, sa mémoire et sa masse sont des lectures différenciées du même état APP--TRIT--TPK. Le transport d'un parcours peut modifier les événements et la mémoire en conservant la valeur du lecteur massique. Pour le déterminer, on utilise le caractère complet de \eqref{md:hib:caracter}, ainsi que l'exposant de retour.

\subsection{Séparation et restitution des deux historiques centraux}
Les registres exprimés des deux orientations sont
\begin{equation}
 e_{++}=(2,2,5,-4,-3,-2)^2,\qquad
 e_{--}=(4,3,2,-2,-2,-5)^2,
 \label{md:r27:historias}
\end{equation}
où l'exposant $2$ indique la concaténation de deux blocs. Ils partagent un total nul et le même histogramme absolu. Pour un registre $e\in\mathbb Z^{12}$, soient
\begin{equation}
 \begin{gathered}
 p_i=e_i+e_{i+6},\quad \mathfrak a_i=e_i-e_{i+6},\\
 s_C=\sum_{i=0}^{5}p_i,\quad M_i=p_i-p_5\quad(0\le i<5).
 \end{gathered}
 \label{md:r27:registro-electron}
\end{equation}
\begin{proposition}[Restitution intégrale du registre]
La donnée $(s_C,M_0,\ldots,M_4,\mathfrak a_0,\ldots,\mathfrak a_5)$ détermine le registre par
\begin{equation}
 \begin{gathered}
 p_5=\frac{s_C-\sum_{i=0}^{4}M_i}{6},\quad p_i=p_5+M_i,\\
 e_i=\frac{p_i+\mathfrak a_i}{2},\quad e_{i+6}=\frac{p_i-\mathfrak a_i}{2}.
 \end{gathered}
 \label{md:r27:inversa-electron}
\end{equation}
Son image entière est caractérisée par la divisibilité par six du premier numérateur et la congruence $p_i\equiv \mathfrak a_i\pmod2$ pour chaque $i$.
\end{proposition}
\begin{proof}
Additionner $p_i=p_5+M_i$ pour les cinq premiers indices et ajouter $p_5$ donne $s_C=6p_5+\sum M_i$. Additionner et soustraire les deux définitions de $p_i,\mathfrak a_i$ donne les formules restantes. Les conditions indiquées sont nécessaires pour obtenir des entiers et suffisantes pour les reconstruire ; la substitution restitue toutes les données originales.
\end{proof}
Pour \eqref{md:r27:historias}, $\mathfrak a=0$ et $M_{++}=(8,8,14,-4,-2)$, tandis que $M_{--}=(18,16,14,6,6)$. La restitution distingue les deux historiques bien que certains lecteurs scalaires partagent leur valeur. Les symboles $\mathfrak a_i$ de cette décomposition sont des différences d'événements et se distinguent du caractère aréal $a_0$ par leur domaine.

\subsection{Forme spectrale du lecteur de retour}
Pour $\tau\in\mathbb R^{12}$, avec des indices cycliques, nous définissons
\begin{equation}
 C_s=\sum_{j=0}^{11}\tau_j\tau_{j+s},\quad
 T_k=\sum_{j=0}^{11}\tau_j e^{-2\pi i kj/12},\quad
 \theta_k=2\pi k/12.
 \label{md:r27:fourier}
\end{equation}
La représentation par caractères s'applique après la production du registre. Le lecteur de retour conserve $\Omega=-2C_1-4C_2-6C_3$ et $P_6=C_6/2$.
\begin{proposition}[Décomposition harmonique exacte]
\begin{equation}
 \begin{split}
 C_s&=\frac1{12}\sum_{k=0}^{11}|T_k|^2\cos(s\theta_k),\\
 \Omega&=\frac1{12}\sum_{k=0}^{11}|T_k|^2
 [-2\cos\theta_k-4\cos2\theta_k-6\cos3\theta_k],\\
 P_6&=\frac1{24}\sum_{k=0}^{11}|T_k|^2(-1)^k.
 \end{split}
 \label{md:r27:espectro}
\end{equation}
\end{proposition}
\begin{proof}
Développer $|T_k|^2$ produit $\sum_{j,l}\tau_j\tau_l e^{-i\theta_k(j-l)}$. La somme $\sum_{k=0}^{11}e^{i\theta_km}$ vaut douze si $12$ divise $m$ et zéro sinon. Appliquer cette identité au déplacement $s$ restitue $C_s$. Comme $\tau$ est réel, les parties imaginaires s'annulent par paires $k,12-k$ et il reste le cosinus. Les deux autres formules résultent d'une combinaison linéaire et de $\cos(6\theta_k)=(-1)^k$.
\end{proof}
Pour $\tau_e=(+++---)^2$, seules subsistent les puissances $|T_2|^2=64$, $|T_6|^2=16$, $|T_{10}|^2=64$. Leurs poids dans $\Omega$ sont $7,4,7$, de sorte que $\Omega=80$ et $P_6=6$. On restitue ainsi la partie harmonique de $\kappa_e=80-54\alpha+6\Delta_4$ dans l'équation électronique intégrale. Le registre alterné $(+-)^6$, de même histogramme, possède $|T_6|^2=144$ et $\Omega=48$ : l'ordre est une dépendance effective du lecteur. Les translations cycliques conservent $|T_k|^2$, même si elles modifient la phase de $T_k$ ; cette lecture spectrale préserve les corrélations, et non l'historique entier.

\subsection{Critère de masse transportée et lectures associées}
Soit $\mathcal H_{\mathcal F}$ la fibre d'un domaine fini de parcours et $\{|\Gamma\rangle\}$ sa base. On fixe l'un des lecteurs de retour $r=D$ ou $r=\Omega$ et une coordinatisation commune avec $R_{\rm act}>0$. La loi des masses de \eqref{md:hib:factor-ruta} détermine
\begin{equation}
 \begin{gathered}
 \frac{m_\Gamma^\beta}{m_e^\beta}
 =\chi_{\rm ref}(\sigma_\Gamma-\sigma_e,\eta_\Gamma-\eta_e)
 R_{\rm act}^{\kappa_\Gamma-\kappa_e},\\
 \Lambda_\Gamma=\ln(m_\Gamma^\beta/m_e^\beta).
 \end{gathered}
 \label{md:r27:masa-relativa}
\end{equation}
Le logarithme agit sur un rapport sans dimension. Dans une fibre non scalaire, on conserve l'opérateur et sa réalisation diagonale déclarée ; le domaine fini peut être une orbite ou une multisection, sans sélectionner une cellule arbitraire.

\begin{theorem}[Conservation du lecteur massique dans un domaine stable]
Soit $F$ un transport TPK agissant bijectivement dans le domaine fixé et $V|\Gamma\rangle=|F\Gamma\rangle$. Pour $M|\Gamma\rangle=m_e^\beta e^{\Lambda_\Gamma}|\Gamma\rangle$, on a
\begin{equation}
 [M,V]|\Gamma\rangle=m_e^\beta
 (e^{\Lambda_{F\Gamma}}-e^{\Lambda_\Gamma})|F\Gamma\rangle.
 \label{md:r27:conmutador-masa}
\end{equation}
Ainsi, $[M,V]=0$ équivaut à $\Lambda_{F\Gamma}=\Lambda_\Gamma$ pour tous les parcours. Les paramètres de la coordinatisation étant fixés, la condition est
\begin{equation}
 \begin{aligned}
 &\delta n_A x+\frac{\delta s_C}{6}y+
 \delta k_{\Delta_4}\Delta_4+\delta b_\zeta\zeta_{270}\\
 &\qquad+\sum_h\delta N_h(q_-^h-q_+^h)
 +\delta\kappa\ln R_{\rm act}=0,
 \end{aligned}
 \label{md:r27:criterio-masa}
\end{equation}
où $\delta$ compare le parcours transporté au parcours initial.
\end{theorem}
\begin{proof}
Appliquer $MV$ et $VM$ à chaque vecteur de la base donne \eqref{md:r27:conmutador-masa}. La positivité de $m_e^\beta$ et l'injectivité de l'exponentielle réelle donnent l'équivalence. Prendre le logarithme de \eqref{md:r27:masa-relativa} et soustraire les caractères des deux parcours donne \eqref{md:r27:criterio-masa}. Dans des tours infinies, on exige pour chaque parcours $\sum_h|N_h(q_-^h-q_+^h)|<\infty$ ; cette convergence absolue légitime la soustraction. La finitude du domaine des parcours conserve les opérateurs bornés utilisés dans cette preuve.
\end{proof}

Le transport est d'abord produit par TPK ; l'égalité détermine ensuite ce qu'il conserve. Les événements individuels et la mémoire peuvent changer et se compenser dans \eqref{md:r27:criterio-masa}. En maintenant la section d'action, l'horloge et la réalisation circulaire de \eqref{md:hib:reciprocidad-radial}, la commutation avec $M$ conserve également les lectures de fréquence propre, de rayon gravitationnel, de longueur de Compton réduite et de poids thermique, qui sont des fonctions du même opérateur positif. Pour l'inverse de la masse, on travaille hors du secteur nul. Dans la réalisation marquée de \eqref{md:r27:area-gravedad}, la quatrième lecture prend la forme
\begin{equation}
 \Xi_\Gamma=\frac{\Omega_\Gamma}{\omega_0}
 =\frac{r_{g,\Gamma}}{L_\alpha}
 =\frac{L_\alpha}{\bar\lambda_\Gamma}
 =\frac{E_\Gamma}{b_*\Theta_{{\rm clk},P}\ln3}.
 \label{md:r27:lecturas-conservadas}
\end{equation}
La preuve consiste à substituer les égalités déjà démontrées dans une même coordinatisation. Ce critère décrit la conservation pendant un transport déclaré ; les sélecteurs physiques d'espèce, de représentation et d'interaction conservent leurs propres conditions. La trajectoire et sa lecture harmonique ajoutent de l'information à la masse scalaire sans remplacer ces conditions.
